\documentclass[11pt]{article}

\usepackage[T1]{fontenc}
\usepackage[utf8]{inputenc}
\usepackage{amsmath,amssymb,amsthm}
\usepackage[margin=1in]{geometry}
\usepackage{microtype}
\usepackage[hidelinks]{hyperref}

\newtheorem{theorem}{Theorem}
\newcommand{\C}{\mathbb{C}}
\newcommand{\Sym}{\operatorname{Sym}}
\newcommand{\mon}{\operatorname{mon}}
\newcommand{\diam}{\operatorname{diam}_{+}}
\newcommand{\slie}{\mathfrak{sl}}

\title{An Exact Certificate for the $(n,k)=(4,3)$\\
Monomial Lie-Additive Diameter}
\author{Xiaojian Zeng\\
\small PBC School of Finance, Tsinghua University, 100190\\
\small \texttt{xiaojianzeng@foxmail.com}}
\date{}

\begin{document}
\maketitle

\begin{abstract}
Let $V=\Sym^3(\C^4)$ and
$U=\Sym^3(\operatorname{span}(e_1,e_2,e_3))$.  For the derived natural
action of $\slie_4(\C)$, we prove that the monomial Lie-additive diameter of
$V$ with respect to $U$ is two.  The upper bound is witnessed by the identity
and a product of three explicitly displayed traceless integer matrices.  In
fixed monomial bases, the corresponding $20\times20$ spanning determinant is
the nonzero integer $8412259151430144$.  This resolves only the fixed cubic
$(n,k)=(4,3)$ cell; the all-$(n,k)$ higher-rank question remains unresolved.
\end{abstract}

\section{Context and statement}

Additive diameters for group representations were introduced by Jezernik and
\v{S}penko, together with a Lie-representation analogue
\cite{JezernikSpenko2025}.  Their later work defines the monomial
Lie-additive diameter used here, proves the optimal formula for
$\slie_3(\C)$ acting on all symmetric powers of $\C^3$, and asks whether the
analogous dimension-counting formula holds for all $n\geq3$ and $k\geq1$
\cite[Section~8, Theorem~8.1 and Question~8.5]{JezernikSpenko2026}.
We determine the single higher-rank specialization $n=4$, $k=3$.  No claim
about the other specializations or about worldwide priority is made.

Put $W=\C^4$ with ordered basis $e_1,e_2,e_3,e_4$, and identify
$V=\Sym^3(W)$ with the homogeneous cubics in these basis vectors.  Set
\[
 U=\Sym^3(\operatorname{span}(e_1,e_2,e_3)).
\]
For $X=(x_{ij})\in\slie_4(\C)$, the derived action is
\begin{equation}\label{eq:action}
 \rho(X)f=\sum_{i,j=1}^4 x_{ij}e_i\frac{\partial f}{\partial e_j}.
\end{equation}
Let $\mon(\rho(\slie_4(\C)))$ consist of the identity and all finite products
$\rho(X_1)\cdots\rho(X_r)$ with $X_i\in\slie_4(\C)$.  The diameter
$\diam^{\slie_4(\C),\mathrm{mon}}(V,U)$ is the least $d$ for which $d$
members $m_i$ of this monoid satisfy $\sum_i m_iU=V$.

\begin{theorem}\label{thm:main}
For the pair above,
\[
 \diam^{\slie_4(\C),\mathrm{mon}}
 \bigl(\Sym^3(\C^4),\Sym^3(\C^3)\bigr)=2.
\]
More precisely, the upper bound is attained by $m_1=I_V$ and
$m_2=\rho(X_1)\rho(X_2)\rho(X_3)$, where the rightmost factor acts first and
\[
\begin{aligned}
X_1&=\begin{pmatrix}
 2& 0& 0&-1\\ 1& 0& 0&-2\\ 1&-1& 1& 0\\-1&-1& 1&-3
\end{pmatrix},&
X_2&=\begin{pmatrix}
 2& 1& 2& 2\\ 0&-2& 2&-1\\-2&-1& 2&-1\\ 2& 1&-2&-2
\end{pmatrix},\\[4pt]
X_3&=\begin{pmatrix}
 1& 1& 0& 0\\ 1& 1& 0& 1\\ 2& 1&-2& 1\\-2&-2&-1& 0
\end{pmatrix}.
\end{aligned}
\]
\end{theorem}

\section{The exact certificate}

Write a cubic monomial as $e^a=e_1^{a_1}e_2^{a_2}e_3^{a_3}e_4^{a_4}$.
Order the exponent tuples for a basis $\mathcal B_U$ of $U$ as
\begin{align*}
 &(0,0,3,0),(0,1,2,0),(0,2,1,0),(0,3,0,0),(1,0,2,0),\\
 &(1,1,1,0),(1,2,0,0),(2,0,1,0),(2,1,0,0),(3,0,0,0),
\end{align*}
and order a complementary monomial basis $\mathcal B_Q$ by
\begin{align*}
 &(0,0,0,3),(0,0,1,2),(0,0,2,1),(0,1,0,2),(0,1,1,1),\\
 &(0,2,0,1),(1,0,0,2),(1,0,1,1),(1,1,0,1),(2,0,0,1).
\end{align*}
Thus $\mathcal B_U$ is a basis of $U$, while the images of
$\mathcal B_Q$ form a basis of $V/U$.

Let $\pi:V\to V/U$ be the quotient map.  Applying~\eqref{eq:action} first
with $X_3$, then $X_2$, then $X_1$, gives the following coefficient matrix
$Q$ for $\pi m_2|_U$ in the two displayed orders:
\[
\begin{pmatrix}
  12&   2&-10&  12&  0&-11& 16&-12& 20& 24\\
  24& -46& 22&  -6&-38& 56&-18& 90&-40&-72\\
 -54&   5&-24& -12& 77&-37& -4&-84& 20& 72\\
 -12&  -4&-12&   6& -4&  9& 40& 20& 28&-36\\
-102& -10&-30&-114& 34&  4&-46&-38& -8& 72\\
   0& -18&-24& -87&  0&  3& -5&-16& 24& 12\\
 -12& -10& 12&   6& -4&  3& 12& -2& 54&138\\
-108&   6& 14&   6& 12& -9&-20& 14&-36&-96\\
   0& -16& -6&  24&-26&-17& -2&-18& -2& 18\\
   0&   4& -2&  -6&-28&-13& -6& -6& -9&-21
\end{pmatrix}.
\]
Exact integer elimination gives
\begin{equation}\label{eq:det}
 \det Q=8412259151430144\neq0.
\end{equation}

\begin{proof}[Proof of Theorem~\ref{thm:main}]
The traces of $X_1,X_2,X_3$ are respectively
$2+0+1-3$, $2-2+2-2$, and $1+1-2+0$, so every factor lies in
$\slie_4(\C)$.  Hence $m_1$ and $m_2$ belong to the required monoid.

Let $u_1,\ldots,u_{10}$ be the monomials in $\mathcal B_U$, and use the
concatenated ordered basis $(\mathcal B_U,\mathcal B_Q)$ of $V$.  The
$20\times20$ matrix whose columns are
\[
 u_1,\ldots,u_{10},m_2u_1,\ldots,m_2u_{10}
\]
has block form
\[
 \begin{pmatrix}I_{10}&A\\0&Q\end{pmatrix}
\]
for an integer matrix $A$.  Its exact determinant is therefore the integer
in~\eqref{eq:det}.  The twenty columns span $V$, proving
$U+m_2U=V$.

Finally,
$\dim V=\binom{6}{3}=20$ and $\dim U=\binom{5}{3}=10$.  The image of $U$
under any single linear operator has dimension at most ten, so one monomial
translate cannot span $V$.  This gives the matching lower bound.
\end{proof}

\section{Verification and scope}

The displayed quotient matrix and the block determinant reduce the proof to
integer arithmetic.  During preparation, the certificate was recomputed in
two ways: one implementation built the three $20\times20$ representation
matrices and used fraction-free elimination; a second applied the derivations
in~\eqref{eq:action} directly to cubics and evaluated the determinant in exact
integer arithmetic.  Both returned the value in~\eqref{eq:det}; no
floating-point rank decision was used.  This separation of a human-readable
reduction from exact and independently reconstructed computations is also
used in recent finite-certificate work, for example the integer-predicate and
independent-verifier architecture of Zeng, Liu, Ratnavelu, and Huat
\cite[Sections~4 and~6]{ZengEtAl2026Certificates}.

The certificate proves the original fixed $(4,3)$ existence statement and
its dimension lower bound.  It supplies no uniform construction for other
$n$ or $k$ and therefore does not resolve Question~8.5 of
\cite{JezernikSpenko2026} in general.

\paragraph{Computational and AI assistance.}
OpenAI Codex assisted with the search for the explicit matrices, exact-check
implementations, literature screening, adversarial audit organization, and
initial manuscript preparation.  One verifier used pure-Python
fraction-free elimination and another used SymPy 1.14.0 for exact symbolic
differentiation and integer determinants.  Language-model output and
floating-point computation are not used as mathematical evidence; the proof
rests on the displayed matrices, action formula, quotient matrix, and exact
nonzero determinant.  No proof assistant was used.

\bibliographystyle{abbrv}
\bibliography{references}

\end{document}
